% Source: MQ12b_v1.html
\documentclass[11pt]{article}
\usepackage[margin=1in]{geometry}
\usepackage{amsmath}
\usepackage{amssymb}
\usepackage{siunitx}
\usepackage{enumitem}

\title{MQ12b: Fluid Mechanics}
\author{}
\date{}

\begin{document}
\maketitle

\section*{MQ12b: Fluid Mechanics}

\begin{enumerate}[leftmargin=*,label=\textbf{Question \arabic*.}]
\item \hfill \textit{1 pts}

An object of mass 3.6 kg and volume 0.6 m$^{3}$ is completely submerged in a tank of fluid with density 632 kg/m$^{3}$. If the acceleration due to gravity is 3.3 m/s$^{2}$at this location (another planet!), calculate the buoyant force \emph{B} on the object.



\emph{B} = \rule{2cm}{0.4pt} N

\item \hfill \textit{1 pts}

A solid wooden raft having density 540 kg/m$^{3}$ and dimensions 4.6 m by 3 m by 0.5 m, where the latter is the thickness, is floating in water. How much mass can be loaded on the raft before 97\% of its volume is submerged?



\emph{m }= \rule{2cm}{0.4pt} kg

\item \hfill \textit{1 pts}

Water flows at 1.4 m/s through a hose with an internal diameter of 1.24 cm, exiting through a nozzle at 4.2 m/s. Calculate the nozzle's diameter in cm.



\emph{D = \rule{2cm}{0.4pt} }cm

\item \hfill \textit{1 pts}

A pipe of cross sectional area 2.43 m$^{2}$ narrows while splitting into three pipes, each of cross-sectional area 0.8 m$^{2}$. If the fluid is traveling at 2.3 m/s in the larger pipe, calculate the fluid speed in each of the narrower pipes.



\emph{v }= \rule{2cm}{0.4pt} m/s

\item \hfill \textit{1 pts}

A chamber 51 m beneath the Earth contains hot water with a density of 1000 kg/m$^{3}$ at a pressure of 729,924 Pa. A narrowing channel leads upwards to a circular vent at the surface having negligible area compared to that of the chamber. At what speed is hot water ejected from the vent? Atmospheric pressure is 1.01 x 10$^{5}$ Pa



\emph{v }= \rule{2cm}{0.4pt} m/s

\item \hfill \textit{1 pts}

A horizontal pipe of radius 0.24 m carries water at 2.4 m/s and pressure of 201 kPa. If the pipe narrows to a radius of 0.13 m, calculate the pressure \emph{p} in kPa in the narrow pipe. The density of water is 1000 kg/m$^{3}$.



\emph{p} = \rule{2cm}{0.4pt} kPa

\end{enumerate}

\end{document}
